\begin{foldable}
  We conduct a wide range of experiments to show our improvements over current SOTA few-shot KD baselines.
  We also include cross-architecture KD, visualization, ablation studies, and comparison with data-free KD methods.
  We provide more details in the \hyperref[sec:supp]{Supplementary}.
\end{foldable}

\subsection{Architectures and datasets}
\begin{foldable}
  We evaluate our method across seven benchmark image datasets: MNIST~\cite{lecun2002gradient}, FMNIST~\cite{xiao2017fashion}, SVHN~\cite{netzer2011reading}, CIFAR10, CIFAR100~\cite{krizhevsky2009learning}, Tiny-ImageNet~\cite{le2015tiny}, and Imagenette~\cite{howard2020imagenette}.
  These datasets cover different difficulty levels, including: (1) \textit{Simple} datasets with MNIST, FMNIST, SVHN, and CIFAR10 (10 classes) and (2) \textit{Complex} datasets with CIFAR100 (100 classes), Tiny-ImageNet (200 classes), and Imagenette (high-resolution at 224\texttimes 224).
\end{foldable}

\begin{foldable}
  We consider several popular network architectures for distillation, namely LeNet5~\cite{lecun2002gradient}, AlexNet~\cite{krizhevsky2012imagenet}, and ResNet~\cite{he2016deep}, and VGG~\cite{simonyan2014very}.
  These architectures and datasets are widely used to evaluate few-shot KD methods~\cite{kimura2018few,wang2020neural,nguyen2022black}.
\end{foldable}

\subsection{Baselines}
\begin{foldable}
  We compare our method DivBFKD with the following baselines:
  \begin{itemize}
    \item \textit{Student-Full}: The student is trained on the whole teacher's training set.
    \item \textit{Student-Alone}: The student is trained on the few-shot set $\mathcal{D}$.
    \item {
        \textit{Standard-KD}~\cite{hinton2015distilling}: The student is trained with the standard KD loss (Eq.~\ref{eq:StandardKD}) on $\mathcal{D}$.
        We select temperature $\lambda=0.9$, a standard value used in KD methods~\cite{hinton2015distilling,tian2020contrastive,yuan2020revisiting,ma2021undistillable}.
      }
    \item \textit{Few-shot KD} methods: We compare with two white-box methods, FSKD~\cite{kimura2018few} and WaGe~\cite{kong2020learning}, and two black-box methods, BBKD~\cite{wang2020neural} and FS-BBT~\cite{nguyen2022black}.
  \end{itemize}
\end{foldable}

\begin{foldable}
  For fair comparisons, we use the same teacher-student network architecture and the same numbers of real and synthetic images $N$ and $M$ as in other few-shot KD methods~\cite{kimura2018few,wang2020neural,nguyen2022black}.
  Their accuracy numbers are obtained from~\cite{nguyen2022black}\footnote{This is possible because we use benchmark datasets with fixed train and test splits.}.
  For our method, we set the quantile $q = 0.1$ to select high-confidence images across all experiments.
  We repeat each experiment three times with random seeds and report the averaged accuracy with standard deviation on the hold-out test set.
\end{foldable}

\subsection{Distillation performance}

\subsubsection{Simple datasets} \label{sec:experiments-simple-datasets}
\begin{table}[ht]
  \centering
  \begin{threeparttable}
    \caption{
      Classification accuracy (\%) on Simple datasets.
      $N$ and $M$ show the budget of real and synthetic images.
    }
    \label{tab:04-experiments-results_few_classes}
    \begin{tabular}{l l l l l}
      \toprule
      \textbf{Dataset} & \textbf{MNIST} & \textbf{FMNIST} & \textbf{SVHN} & \textbf{CIFAR10} \\
      $(N,M)$ & (2\,K, 24\,K) & (2\,K, 48\,K) & (2\,K, 40\,K) & (2\,K, 40\,K) \\
      \midrule
      Teacher                   & 99.28                   & 90.90                   & 96.16                   & 90.05 \\
      Student-Full              & \meanstd{98.91}{0.08}   & \meanstd{88.68}{0.53}   & \meanstd{95.61}{0.07}   & \meanstd{87.58}{0.36} \\
      \midrule

      \multicolumn{5}{c}{\textbf{Smaller Architecture}} \\
      \midrule
      Student-Alone             & \meanstd{95.14}{0.19}   & \meanstd{79.24}{1.05}   & \meanstd{85.68}{0.46}   & \meanstd{59.02}{0.54} \\
      Standard-KD\blackbox      & \meanstd{95.36}{0.11}   & \meanstd{80.76}{0.37}   & \meanstd{89.11}{0.22}   & \meanstd{59.77}{0.96} \\
      FSKD\whitebox             & 80.43\dg                & 68.64\dg                & $-$                     & 40.58\dg \\
      WaGe\whitebox             & $-$                     & 85.18\dg                & $-$                     & 73.08\dg \\
      FS-BBT\blackbox           & 98.42\dg                & 84.73\dg                & $-$                     & 74.10\dg \\
      \textbf{DivBFKD}\blackbox & \meanstdbf{98.58}{0.06} & \meanstdbf{86.50}{0.18} & \meanstdbf{94.47}{0.12} & \meanstdbf{76.97}{0.56} \\
      \midrule

      \multicolumn{5}{c}{\textbf{Same Architecture}} \\
      \midrule
      BBKD\blackbox             & 98.74\dg                & 80.90\dg                & $-$                     & 74.60\dg \\
      FS-BBT\blackbox           & 98.91\dg                & 86.53\dg                & $-$                     & 76.17\dg \\
      \textbf{DivBFKD}\blackbox & \meanstdbf{98.97}{0.04} & \meanstdbf{87.65}{0.37} & \meanstdbf{94.61}{0.08} & \meanstdbf{78.07}{0.26} \\
      \bottomrule
    \end{tabular}
    \begin{tablenotes}[flushleft] \footnotesize
    \item {
        [(W)] white-box; [(B)] black-box; [$-$] not reported; [$\dagger$] reported from respective paper; [\textbf{bold}] best results.
      }
    \end{tablenotes}
  \end{threeparttable}
\end{table}

\begin{foldable}
  Following~\cite{nguyen2022black}, we conduct distillation in two settings:
  \begin{enumerate}
    \item \textit{Smaller} architecture: The student has a smaller architecture than the teacher: LeNet5 to LeNet5-Half on MNIST, FMNIST; and AlexNet to AlexNet-Half on SVHN, CIFAR10.\footnote{The \textit{-Half} version has half the channels in each layer.}
    \item \textit{Same} architecture: The student has an identical architecture to the teacher.
  \end{enumerate}
\end{foldable}

\begin{foldable}
  We report the results in Table~\ref{tab:04-experiments-results_few_classes} and make the following observations:
  \begin{itemize}
    \item Our method DivBFKD improves by around 3-6\% over Standard-KD on MNIST, FMNIST, and SVHN and by significantly 17\% on CIFAR10.
    \item {
        Compared to other few-shot KD methods, our DivBFKD also excels in both smaller and same architecture settings.
        It improves upon second-best FS-BBT by around 2\% on FMNIST and CIFAR10.
      }
    \item {
        The performance gap between our DivBFKD and its upper bound Student-Full is very close on MNIST (0.33\%), FMNIST (2.18\%), and SVHN (1.14\%).
        We emphasize that Student-Full uses the teacher's training set ($\ge$ 60\,K real images), much larger than the 2\,K real images we started with.
      }
  \end{itemize}
\end{foldable}

\subsubsection{Complex datasets} \label{sec:experiments-complex-datasets}
\begin{foldable}
  Similar to~\cite{wang2020neural,nguyen2022black,kong2020learning}, we distill from ResNet32 to ResNet20 on CIFAR100 and Tiny-ImageNet (Table~\ref{tab:ResultsHundredClasses}).
  With limited data, Student-Alone loses roughly half of its accuracy compared to Student-Full.
  As a naive solution, Standard-KD improves over Student-Alone 10\% on CIFAR100 and 14\% on Tiny-ImageNet.
  Ultimately, our method DivBFKD is the best method, exceeding the second-best FS-BBT by 3\% on CIFAR100 and 1\% on Tiny-ImageNet.
\end{foldable}

\begin{foldable}
  We further evaluate our method on high-resolution images with Imagenette.
  On 2\,K real images, Student-Alone achieves 76.80\% while Standard-KD marginally improves to 76.94\%.
  DivBFKD further improves upon this result by 10\%, approximating Student-Full.
\end{foldable}

\begin{table}[ht]
  \centering
  \begin{threeparttable}
    \caption{
      Classification accuracy (\%) on Complex datasets.
      $N$ and $M$ show the budget of real and synthetic images.
    }
    \label{tab:ResultsHundredClasses}
    \begin{tabular}{l l l l}
      \toprule
      \textbf{Dataset} & \textbf{CIFAR100} & \textbf{Tiny-ImageNet} & \textbf{Imagenette} \\
      $(N,M)$                   & (5\,K, 40\,K)           & (10\,K, 50\,K)        & (2\,K, 40\,K) \\
      \midrule
      Teacher                   & 71.41                   & 56.27                 & 90.47 \\
      Student-Full              & \meanstd{68.28}{0.20}   & \meanstd{51.51}{0.31} & \meanstd{89.89}{0.41} \\
      \midrule
      Student-Alone             & \meanstd{35.71}{0.39}   & \meanstd{26.59}{0.25} & \meanstd{76.80}{0.78} \\
      Standard-KD\blackbox      & \meanstd{45.76}{0.48}   & \meanstd{40.42}{0.19} & \meanstd{76.94}{0.14} \\
      WaGe\whitebox             & 20.32\dg                & $-$                   & $-$ \\
      BBKD\blackbox             & 53.41\dg                & 40.01\dg              & $-$ \\
      FS-BBT\blackbox           & 56.28\dg                & 43.29\dg              & $-$ \\
      \textbf{DivBFKD}\blackbox & \meanstdbf{59.91}{0.31} & \meanstdbf{44.25}{0.38} & \meanstdbf{86.48}{0.72} \\
      \bottomrule
    \end{tabular}
    \begin{tablenotes}[flushleft] \footnotesize
    \item[] {
        [(W)] white-box; [(B)] black-box; [$-$] not reported; [$\dagger$] reported from respective paper; [\textbf{bold}] best results.
      }
    \end{tablenotes}
  \end{threeparttable}
\end{table}

\subsubsection{Cross-architecture distillation} \label{sec:experiments-crossarch}
\begin{foldable}
  In previous subsections, we distill from the teacher to a student networks of similar architectures to have a consistent setting with other few-shot KD methods (\eg, LeNet5 to LeNet5-Half, AlexNet to AlexNet-Half, and ResNet32 to ResNet20).
  This section further explores the cases where the teacher and student architectures belong to different families: ResNet~\cite{he2016deep}, VGG~\cite{simonyan2014very}, and AlexNet~\cite{krizhevsky2012imagenet}.
  We adopt the same setting from the standard experiment on CIFAR10 from \S\ref{sec:experiments-complex-datasets}.
  Table~\ref{tab:Cross-architecture} shows the performance of our method DivBFKD with three teachers ResNet32, VGG16, and AlexNet combined with three students ResNet20, VGG11, and AlexNet-Half on CIFAR10.
\end{foldable}

\begin{table}[ht]
  \caption{Student accuracy (\%) of cross-architecture distillation on CIFAR10.}
  \label{tab:Cross-architecture}
  \centering
  \begin{threeparttable}
    \begin{tabular}{c c|c c c c c}
      \toprule
      \multicolumn{2}{c|}{\textbf{Architecture}} & \multicolumn{5}{c}{\textbf{Accuracy}}\\
      Teacher & Student & Teacher & S-Full & S-Alone & Standard-KD & DivBFKD \\
      \midrule
      ResNet32 & ResNet20     & 93.34 & & & \meanstd{65.86}{1.90} & \meanstd{84.05}{0.10} \\
      VGG16    & ResNet20     & 91.30 & \meanstd{92.34}{0.06} & \meanstd{64.52}{1.60} & \meanstd{64.76}{0.89} & \meanstd{80.69}{0.48} \\
      AlexNet  & ResNet20     & 90.05 & & & \meanstd{65.29}{0.89} & \meanstd{79.57}{0.16} \\
      \midrule
      ResNet32 & VGG11        & 93.34 & & & \meanstd{53.78}{0.78} & \meanstd{69.42}{0.83} \\
      VGG16    & VGG11        & 91.30 & \meanstd{89.12}{0.17} & \meanstd{53.34}{0.57} & \meanstd{53.80}{1.21} & \meanstd{75.41}{0.06} \\
      AlexNet  & VGG11        & 90.05 & & & \meanstd{53.51}{0.14} & \meanstd{74.60}{0.29} \\
      \midrule
      ResNet32 & AlexNet-Half & 93.34 & & & \meanstd{60.36}{0.82} & \meanstd{74.70}{0.31} \\
      VGG16    & AlexNet-Half & 91.30 & \meanstd{87.58}{0.36} & \meanstd{59.02}{0.54} & \meanstd{60.47}{0.10} & \meanstd{76.16}{0.57} \\
      AlexNet  & AlexNet-Half & 90.05 & & & \meanstd{59.77}{0.96} & \meanstd{76.97}{0.56} \\
      \bottomrule
    \end{tabular}
    \begin{tablenotes}[flushleft] \footnotesize
    \item {Same teacher-student pairs share the same Student-Full and Student-Alone accuracy.}
    \end{tablenotes}
  \end{threeparttable}
\end{table}

\begin{foldable}
  We have the following observations:
  \begin{itemize}
    \item Among three architecture families, ResNet is the best one with a 93.34\% Teacher, 92.34\% Student-Full, and 64.52\% Student-Alone.
    \item Standard-KD is helpful and it is always better than Student-Alone in all teacher-student combinations.
    \item {
        VGG is better than AlexNet when the teacher and the student are trained on the whole training set.
        However, VGG student is worse than AlexNet student when it is trained on the few-shot dataset.
        Namely, Student-Alone of VGG11 has 53.34\% accuracy while Student-Alone of AlexNet-Half has 59.02\% accuracy.
      }
    \item Our method DivBFKD remarkably improves accuracy over Standard-KD by 14--22\% on all teacher-student combinations.
    \item Distillation is often most effective when the teacher's and student's architecture belong to the same family.
  \end{itemize}
  Overall, the robust performance across multiple data scales and architectures suggests that our method successfully tackles the black-box few-shot KD setting.
\end{foldable}

\subsection{Diversity of our synthetic images}

\subsubsection{Comparison with the teacher's training images.}

\begin{foldable}
  As we mentioned earlier, our goal is to generate synthetic images close to the teacher's training images.
  To prove this argument, we embed and visualize the teacher's training images (they are \textit{unknown} and not used in the KD process) and our synthetic images of CIFAR10 in Fig.~\ref{fig:Distributions-of-CIFAR10}.
  It shows that our synthetic images belong to the distribution of the teacher's training images.
  With the teacher's supervision, our synthetic images show their strong resemblance to the teacher's training images.
\end{foldable}

\begin{foldable}
  For quantitative results, we compute the \textit{Coverage} metric~\cite{alaa2022faithful} to measure the fraction of teacher's training samples whose neighborhoods contain at least one synthetic sample.
  Our synthetic images achieve 0.42, compared to 0.18 of the few-shot images and 0.29 of the synthetic images generated from the standard WGAN.
  This shows that our synthetic images resemble many variability of the teacher's training images, thus improving image diversity.
\end{foldable}

\begin{figure}[ht]
  \centering
  \includegraphics[width=1\columnwidth]{./figures/04-experiments-viz-embed2d.png}
  \caption{Embeddings of teacher's training images, few-shot images, and our synthetic images on CIFAR10.}
  \label{fig:Distributions-of-CIFAR10}
\end{figure}

\begin{figure}[ht]
  \centering
  \begin{subfigure}[t]{0.475\columnwidth}
    \centering
    \includegraphics[width=\linewidth]{./figures/04-experiments-viz-real.pdf}
    \caption{Real images}
  \end{subfigure}
  \hfill
  \begin{subfigure}[t]{0.475\columnwidth}
    \centering
    \includegraphics[width=\linewidth]{./figures/04-experiments-viz-synth.pdf}
    \caption{Synthetic images}
  \end{subfigure}
  \caption{Visualization of real and synthetic images along with their predictive labels and confidence-scores provided by the teacher on CIFAR10.}
  \label{fig:Visualization}
\end{figure}

\subsubsection{Quality and quantitative measurements}
\begin{foldable}
  We randomly visualize 24 real images and 24 synthetic images generated by our WGAN in Fig.~\ref{fig:Visualization}.
  Primarily, our goal is not to generate visually aesthetic images but to generate diverse images that the teacher confidently predicts their labels.
  As a side note, image quality is also not a strict requirement in KD methods relied on synthetic images~\cite{chen2019data,wang2020neural,wang2021zero}.
  However, our synthetic images in Fig.~\ref{fig:Visualization}(b) are able to show convincing properties of their real counterparts.
\end{foldable}

\begin{foldable}
  \textbf{Inception score (IS) and Fr\'{e}chet inception distance (FID).} To quantify the diversity and quality of synthetic images, we compute the IS~\cite{salimans2016improved} (higher is better) and FID~\cite{heusel2017gans} (lower is better) of our synthetic images in Table~\ref{tab:IS-FID}.
  We compare with the standard WGAN and FS-BBT (the second-best baseline), where our method achieves the best scores.
  These results confirm that our method generate synthetic images with high diversity and quality, and explain their effectiveness for the distillation step.
\end{foldable}

\begin{table}[ht]
  \centering
  \begin{threeparttable}
    \caption{IS and FID of synthetic images on CIFAR10.}
    \label{tab:IS-FID}
    \begin{tabular}{l c c}
      \toprule
      \textbf{Method} & \textbf{IS ($\uparrow$)} & \textbf{FID ($\downarrow$)} \\
      \midrule
      CVAE (from~\cite{nguyen2022black}) & 3.51          & 19.91 \\
      Standard WGAN                      & 3.99          & 14.79 \\
      WGAN with AT (Ours)                & \textbf{4.42} & \textbf{13.96} \\
      \bottomrule
    \end{tabular}
    \begin{tablenotes}[flushleft] \footnotesize
    \item {[AT] adaptive thresholds.}
    \end{tablenotes}
  \end{threeparttable}
\end{table}

\subsection{Ablation studies}

\subsubsection{Impacts of different components}
\begin{foldable}

  As described in \S\ref{sec:03-framework}, the core components of our method include WGAN and adaptive thresholds in the generation phase, and class balancing in the distillation phase.
  Adding both latter components to WGAN completes our method.
  Table~\ref{tab:Ablation-study-Components} breaks down the contribution of each component, where both adaptive thresholds (AT) and class balancing (CB) improve our accuracy when applied either independently or jointly.
\end{foldable}

\begin{table}[ht]
  \centering
  \begin{threeparttable}
    \caption{Student accuracy (\%) with incremental components of DivBFKD.}
    \label{tab:Ablation-study-Components}
    \begin{tabular}{l c c}
      \toprule
      \textbf{Components}                 & \textbf{CIFAR10}        & \textbf{Tiny-ImageNet} \\
      \midrule
      WGAN~\cite{arjovsky2017wasserstein} & \meanstd{75.59}{0.69}   & \meanstd{42.70}{0.25} \\
      WGAN with AT                        & \meanstd{76.35}{0.21}   & \meanstd{43.25}{0.71} \\
      WGAN with CB                        & \meanstd{76.88}{0.45}   & \meanstd{43.75}{0.40} \\
      \textbf{DivBFKD}                    & \meanstdbf{76.97}{0.56} & \meanstdbf{44.25}{0.38} \\
      \bottomrule
    \end{tabular}
    \begin{tablenotes}[flushleft] \footnotesize
    \item {[AT] adaptive thresholds; [CB] class balancing.}
    \end{tablenotes}

  \end{threeparttable}
\end{table}

\subsubsection{Impacts of the numbers of real and synthetic images}
\begin{foldable}
  We investigate our performance with varying amounts of data on CIFAR10.
  Fig.~\ref{fig:04-experiments-abla_data}(a) shows that both Standard-KD and our DivBFKD improve student accuracy with more real images, as expected.
  At $N$ = 5\,K, our DivBFKD achieves 79.93\%, while Standard-KD only achieves 72.66\%.
  Conversely, at $N$ = 250, Standard-KD drops abruptly to 33.33\%, while our method still has an acceptable 66.76\%.
  This trend proves that our method can still distill robustly in extremely limited data cases.
  Besides, we also expect that our performance should improve with more synthetic images.
  This agrees with Fig.~\ref{fig:04-experiments-abla_data}(b) where our DivBFKD clearly outperforms Standard-KD and improves with larger $M$.
\end{foldable}

\begin{figure}[ht]
  \centering
  \includegraphics[width=0.485\columnwidth]{./figures/04-experiments-ablation-n.pdf}\hspace{0.2cm}\includegraphics[width=0.485\columnwidth]{./figures/04-experiments-ablation-m.pdf}
  (a)\hspace{0.49\columnwidth}(b)
  \caption{Accuracy versus the budget of real ($N$) and synthetic images ($M$) on CIFAR10.}
  \label{fig:04-experiments-abla_data}
\end{figure}

\subsubsection{Impact of the quantile}
\begin{foldable}
  In our method, the quantile $q$ is principally used to compute adaptive thresholds to determine high-confidence images.
  Contrarily, while too small $q$ might introduce counter-productive noise, too large $q$ does not allow enough high-confidence images to improve diversity.
  We achieve a relatively stable accuracy of 76.22--76.97\% with $q\in[0.03,0.10]$ (Fig.~\ref{fig:Results-Hyperparameter-tuning}) and recommend this range for a good trade-off between image diversity and training stability.
\end{foldable}

\begin{figure}[ht]
  \centering
  \includegraphics[width=0.75\columnwidth]{./figures/04-experiments-hyperparam-q.pdf}
  \caption{Accuracy versus the quantile $q$ on CIFAR10.}
  \label{fig:Results-Hyperparameter-tuning}
\end{figure}

\subsection{Comparison with data-free KD methods} \label{ssec:04-experiments-comparison_zero_shot}
\begin{foldable}
  We also compare our method with several popular data-free (zero-shot) KD methods, including four white-box methods: Meta-KD~\cite{lopes2017data}, ZSKD~\cite{nayak2019zero}, DAFL~\cite{chen2019data}, and DFKD~\cite{wang2021data}, and two black-box methods ZSDB3KD~\cite{wang2021zero} and IDEAL~\cite{zhang2022ideal}.
  From Table~\ref{tab:04-experiments-comparison_data_free}, our method DivBFKD can successfully take advantage of a small amount of data (2\,K images) and perform much better than data-free methods on FMNIST and CIFAR10, while it is comparable on MNIST.
\end{foldable}

\begin{table}[ht]
  \centering
  \begin{threeparttable}
    \caption{Classification accuracy (\%) comparison with zero-shot KD methods.}
    \label{tab:04-experiments-comparison_data_free}
    \begin{tabular}{l|c|c|c}
      \hline
      \textbf{Dataset} & \textbf{MNIST} & \textbf{FMNIST} & \textbf{CIFAR10}  \\
      \hline
      Meta-KD\whitebox          & 92.47          & $-$            & $-$ \\
      ZSKD\whitebox             & 98.77          & 79.62          & 69.56 \\
      DAFL\whitebox             & 98.20          & $-$            & 66.38 \\
      DFKD\whitebox             & \textbf{99.08} & $-$            & 73.91 \\
      ZSDB3KD\blackbox          & 96.54          & 72.31          & 59.46 \\
      IDEAL\blackbox            & 96.32          & 83.92          & 65.61 \\
      \textbf{DivBFKD}\blackbox & 98.97          & \textbf{87.65} & \textbf{78.07} \\
      \hline
    \end{tabular}
    \begin{tablenotes}[flushleft] \footnotesize
    \item {[(W)] white-box; [(B)] black-box; [$-$] not reported; [\textbf{bold}] best results.}
    \end{tablenotes}
  \end{threeparttable}
\end{table}
